\section{Coverage and order of limits}
\label{app:v107-coverage}
The organizing event and pointwise target are still
\eqref{eq:v76-event} and \eqref{eq:v76-target}. The following account
separates established estimates for positive source classes from the
remaining estimate on their complement.

\paragraph{Positive residence.}
The persistence theorem already estimates every selected physical
sheet with $\ell\ge h_B$ at a previously fixed $h_B>0$. This includes
flat isolated pulses, singular endpoint types and reach-deficient
sheets whenever their total residence is at least $h_B$. Uniformity
of the leading constant in a fixed $h_B$ does not permit replacing
$h_B$ by a collision-dependent number.

\paragraph{Stationary and moving comparison families.}
The older critical-arm and adaptive-reach theorems retain their
proved classes. Theorem~\ref{thm:v107-envelope} additionally applies
to nonstationary entrances, nonstationary exits and weighted
comparison overlaps, provided actual interval inclusion and
comparison capacity hold. Its proof uses only regular-segment
variation and the original volume source. A singular endpoint type
is therefore not excluded by its name, but it must satisfy those
actual input and comparison hypotheses; the regular critical Hessian
formula is not extended to singular endpoints by terminology.

\paragraph{Trace loss versus capacity shortage.}
For each comparison block the trace loss is $E_j=P_j-D_j$ and is
paid in full by \eqref{eq:v107-finite}. The remaining capacity
shortage is $(C_j-\kappa P_j)_+$, not $(C_j-\kappa D_j)_+$.
Theorem~\ref{thm:v107-allocation} proves this distinction even when
$D_j=0$. Formula~\eqref{eq:v107-R} retains both the uncovered
microscopic sheets and the unallocated fractions. It is not a
signed error, a moment-truncated substitute or an omitted face term.

\paragraph{Actual reach and uniformity.}
All comparison intervals lie on actual sheets of the same exact
event. The overlap bound is imposed on their full incident-sheet
supports; disjoint starting sets alone do not establish it.
Proposition~\ref{prop:v107-drifting} is a sufficient verification
on a stated class of arms. Finite-count definability alone supplies
neither its level-width constant nor a budget uniform in $m,R,n,k$.
Consequently the proposition is not a classification of the entire
physical source.

\paragraph{Quantifiers.}
First fix $M,B,\varepsilon(B),h_B,r_B,\Delta_B$ and all other
auxiliary bands. At every finite $m,R,n,k$ form the entire physical
sum and its exact nonnegative decomposition. Next take the roof
essential supremum, the supremum in radius and exact labels, and the
collision limit superior in the order defining $\mathcal H_M$.
Only afterwards may the fixed auxiliary bands be enlarged or a
previously fixed residence threshold be decreased. The outer
parameter $B$ tends to infinity last. The estimate for a complete
shrinking-residence tail is not obtained by exchanging the collision
limit with $h\downarrow0$. The conditional-law conclusions retain
their original scalar-first order.

\paragraph{Historical provenance and present conclusion.}
The submission-state review of revision 103 predates its final
mathematical push; it is not a substantive rejection of that later
source. The substantive report addressed here is the revision-104
report. Revisions 105 and 106 are retained as subsequent author
work, not described as externally certified mathematics. Every
inherited numbered mathematical module is retained. The malformed
paid-arm reference in the active introduction has been corrected;
the historical source snapshot remains byte-preserved.

The new conclusions are Theorems~\ref{thm:v107-envelope},
\ref{thm:v107-allocation} and \ref{thm:v107-original}, with their
stated hypotheses. They directly enlarge the paid positive source
and sharpen its exact remainder. The global estimate
\eqref{eq:v107-final-height} is not proved here. In particular the
raw pointwise theorem, unrestricted conditioned laws and an
independent specialist verification are not certified by the finite
checks or by successful typesetting.
